\section{The trace-73 Cappell--Shaneson candidate}

Let
\begin{equation}\label{eq:A}
A=\begin{pmatrix}
0&269&1240\\
0&41&189\\
1&0&32
\end{pmatrix}.
\end{equation}
This is Iwaki's standard form $X_{41,189,73}$; the triple occurs among the
trace-73 representatives in his table \cite[Table 2]{Iwaki2024}.  Direct
integer arithmetic gives
\begin{equation}\label{eq:determinants}
\det A=1,\qquad \det(A-I)=1.
\end{equation}

For $A\in \operatorname{SL}(3,\Z)$, let $f_A:T^3\to T^3$ be the induced
diffeomorphism, made equal to the identity near a chosen point.  Surgery on
the corresponding circle in the mapping torus gives a
Cappell--Shaneson manifold $\Sigma_A^\varepsilon$, where
$\varepsilon\in\Z/2$ records the framing choice.  The standard criterion says
that $\Sigma_A^\varepsilon$ is a homotopy 4-sphere if and only if
$\det(A-I)=\pm1$; see Iwaki's Proposition 2.1 and the original sources cited
there \cite{Iwaki2024,CappellShaneson1976}.

Consequently, \eqref{eq:determinants} proves that the standard construction
$\Sigma_A^\varepsilon$ is a homotopy 4-sphere.  It does not prove that a
separately generated large Kirby diagram is a diagram of that construction.
That identification is the first and most load-bearing issue below.

Iwaki proves standardness for many infinite families.  The appearance of
$(41,189,73)$ in a representative table is not a standardness or exoticness
theorem for this individual sphere.  Likewise, the trace bound in
Kim--Yamada's work stops below 73 \cite{KimYamada2017}.  These observations
explain the choice of candidate but do not decide its diffeomorphism type.
